\documentclass[10pt,a4paper]{article}
\usepackage[T1]{fontenc}
\usepackage[utf8]{inputenc}
\usepackage{lmodern}
\usepackage[margin=2.2cm]{geometry}
\usepackage{amsmath,amssymb}
\usepackage{booktabs}
\usepackage{upgreek}
\usepackage{enumitem}
\usepackage[colorlinks=true,linkcolor=blue!50!black,urlcolor=blue!50!black]{hyperref}
\newcommand{\ket}[1]{|#1\rangle}
\newcommand{\etaeff}{\eta_{\perp}}
\newcommand{\SI}[2]{\ensuremath{#1\,\mathrm{#2}}}
\newcommand{\more}[2]{\href{report.pdf\##1}{[#2]}}
\setlist{nosep,leftmargin=1.2em}
\title{Radial EIT cooling in a hollow-core fibre: core facts and equations}
\author{}
\date{November 2026}
\begin{document}
\maketitle
\noindent Facts and equations only. Derivations and evidence are in the \href{report.pdf}{full report}; links in brackets go to the section (keep both PDFs in one folder). Status: \textbf{[N]} master-equation numerics; \textbf{[C]} rate chain; \textbf{[M]} analytic model, untested or partly tested. Unless stated: closed $\Lambda$ system, no Zeeman structure, $^{87}$Rb D1, $\lambda_L=\SI{1064}{nm}$.

\section*{Setting}
Coupling in LP$_{01}$, probe in LP$_{11}$, atoms on its node: the carrier vanishes, and the probe couples through its gradient, $\etaeff=G_{11}x_0/a$, $G_{11}=3.49$. Narrow line $\Gamma_n\simeq\Gamma\nu_r/\Delta$; saturation $s=2\Omega_R^2/\Gamma_n^2$. \more{d:soft}{report}

\section*{Control parameters}
\begin{equation*}
  h=\frac{\alpha}{b_p}\frac{\omega_{\rm rec}}{\nu_r}=0.188\,\frac{\zeta}{b_p},\qquad
  \zeta\equiv k_La\,\frac{\omega_{\rm rec}}{\nu_z},\qquad
  \frac{\nu_r}{\omega_{\rm rec}}=\frac{2.11}{\zeta},\qquad
  \etaeff=\frac{G_{11}}{ka}\sqrt{\frac{\zeta}{2.11}}
\end{equation*}
$h$ is the recoil per removed phonon, $\alpha=0.397$. Case A ($a=\ensuremath{15\,\upmu\mathrm{m}}$, $\nu_z=\SI{300}{kHz}$): $\zeta=1.07$. Case B ($a=\ensuremath{3.5\,\upmu\mathrm{m}}$): $\zeta=0.25$. \textbf{The core radius does not fix $h$: $\nu_z$ compensates}, at a cost in power ($\propto\nu_z^2a^2$), purity and excited-state shift. \more{d:zeta}{control parameter}

\section*{Rates}
\begin{itemize}
\item \textbf{[N]} $W$, from $\langle n\rangle$, depends on $N$, with a pole near $n_{\rm ss}\approx0.2$. $W_0^{-1}=\int(P_0^{\rm ss}-P_0)\,dt/P_0^{\rm ss}$ is stable. Case A, $s=1$: $W=1.6\times10^3$, $W_0=2.2\times10^3\,$s$^{-1}$, $n_{\rm ss}=0.15$, $P_0=0.953$. $W_0$ is meaningful only for $P_0\gtrsim0.85$. \more{d:W0}{rates}
\item \textbf{[M, tested on data]} One formula for all three regimes, with $n_c=2(\nu_r/\Gamma_n)^2/s$: \more{d:regimes}{derivation}
\begin{equation*}
  R(n)\simeq R_{\rm sat}\frac{c_{\rm sat}sn}{1+c_{\rm sat}sn}\Big(1+c\frac{n}{n_c}\Big),\qquad c\approx1.3\text{--}1.9
\end{equation*}
The weak limit is the red-sideband rate $ns\Gamma_n/2$; the strong limit is gray molasses. \textbf{[N]} Confirmed: the onset moves with $n_c$ ($\Delta=2$, 3, $5\Gamma$), and at $\Delta<0$ the fast channel heats ($+2\times10^4\,$s$^{-1}$ at $n=200$), leaving a capture range below $n\approx8$. The dip sits at $n\approx0.5\,n_c$ and moves with the onset; not derived.
\end{itemize}

\section*{Fock statistics}
\begin{itemize}
\item \textbf{[N, M]} $P(n)\propto n^{-1/h}e^{-n/n_x}$, $n_x\sim2hn_c$. Measured slope $-2.34$, predicted $-2.48$, at $s=0.1$. In case A at $s=0.1$, $3\times10^{-4}$ of the atoms carry 89\% of $\langle n^2\rangle$. \more{d:tail}{tail}
\item \textbf{[M]} Without a cut-off: variance diverges at $h\ge1/3$, mean at $h\ge1/2$. These are crossovers of width $\delta h\simeq h^2/\ln[2(\Delta/\Gamma)^2]$.
\item \textbf{Prior work:} the same tail law and thresholds are known for Sisyphus cooling in a harmonic trap (Dechant \emph{et al.}, PRL 115, 173006; PRE 94, 022151). Not new. Possibly new: this in ground-state cooling, with the cut-off.
\item \textbf{[T,N]} TE$_{01}$ is an exact, non-degenerate fibre mode (lowest loss of the LP$_{11}$ group, 1.49$\times$ HE$_{11}$ in the tube model) with a node in all field components; LP$_{11}$ polarised along its lobes has an on-axis $E_z=1.4\,\eta_\perp$ (A), $2.9\,\eta_\perp$ (B). Its two transverse components on two $\Lambda$ channels with \emph{separate} final states cool both axes at once: $P_{00}=0.93$--0.94, $W_{00}=0.35$--0.41\,$W_0$, $1.85\times$ switching (2D ME); shared final state: only $x+y$ cools. $^{87}$Rb: $\ket{2,1}$ via $F'=2$, two probe frequencies, $6\times$ coupling imbalance; case A excluded by recoil, case B allowed [estimate]. \more{d:te01}{TE01}
\item \textbf{[N]} Fixed $h$, $\Delta$ and $s$ varied: tail exponent $=1/h$ within 1--5\% wherever the window spans $>1$ decade (up to 1.6); local exponent $\gamma(n)h$ collapses across $h$, but at $s=1$ has structure (dip near $0.3$--$0.45\,n_c$); clean $n^{-1/h}$ needs small $s$, large $\Delta$; band and $N$ converged. \more{d:map}{map}
\item \textbf{[N]} Universality: $h=0.1$--1.2 by two routes ($\nu_r/\omega_{\rm rec}$ or $b_p$), $\Delta=8\Gamma$: the tail exponent collapses on $1/h$ (within 3\% between routes, 4--9\% below $1/h$); $1-P_0$ does not; no sharp threshold. \more{d:map}{map}
\item \textbf{[N]} Map in $(\nu_r/\omega_{\rm rec},b_p)$: the tail measure follows $h$ within about 1--30\% for $h\gtrsim0.4$; $P_0$ and $W_0$ do not. \more{d:map}{map}
\end{itemize}

\section*{Both radial axes}
\begin{itemize}
\item \textbf{[M]} One coherent probe cools one motional mode, $c_xa+c_yb$. LP$_{11}$ cools one axis; \textbf{a scalar OAM beam cools one circular mode}. Options: alternating lobes (slower than $1/\Gamma_n$), rotating pattern, trap anisotropy, two OAM beams, all $\approx W/2$ per axis. TE$_{01}$/TM$_{01}$ may cool both modes, if $\sigma^\pm$ paths are independent. \more{d:2d}{two axes}
\item \textbf{[M, C]} Recoil heats both axes: $h\to2h$ when both are saturated, so
\begin{equation*}
  h<\tfrac12\ \Leftrightarrow\ \zeta<2.66\,b_p\quad\text{(cold state)},\qquad 2h<\tfrac13\ \Leftrightarrow\ \zeta<0.89\,b_p\quad\text{(clean tail)},\qquad n_{\rm core}=\Big(\frac{\Gamma}{4\Delta}\Big)^2\frac{1+2h}{1-2h}
\end{equation*}
\item \textbf{[C]} $P_{00}$ at $s=1$, $3\Gamma$ / $6\Gamma$: case A, $b_p=1/2$: 0.755 / 0.941; case A, $b_p=5/6$: 0.927 / 0.982; $\nu_r/\omega_{\rm rec}=5$, $b_p=5/6$: 0.956 / 0.989. Exact check at $N=6$ per axis: chain $P_{00}$ within 0.6\% for $h\le0.26$; case A, $b_p=5/6$: 0.927 against 0.931. \textbf{[N]} Alternating stages must last $\tau\gtrsim10/\Gamma_n$ ($\approx\SI{0.7}{ms}$ in case A).
\item \textbf{[C]} Alternating the LP$_{11}$ orientation needs stages $\tau\gtrsim5/\Gamma_n$; faster switching collapses $P_{00}$ (0.19 at $\tau\Gamma_n=0.3$). The 2D rate peaks at $W_{00}\approx0.2\,W_0$ for $\tau\Gamma_n\approx5$--15 (0.3--1\,ms in case A), falling for longer stages; $W_{00}P_{00}\approx0.17$--$0.20$ of $W_0P_0$. Exact 2D master equation, converged in $N$ (6, 8, 10). \more{d:2d}{2D}
\end{itemize}

\section*{Preparation time}
\begin{itemize}
\item \textbf{[N]} One axis, $s=1$, time to $P_0\ge0.9$ from $\bar n_0=30$: 20.8\,ms (case A, $b_p=1/2$), 12.7\,ms ($b_p=5/6$), 3.5\,ms (case B). The GRM is 2.3$\times$ and 1.7$\times$ too slow when hot, and cooling time is sublinear in temperature. \more{d:prep}{preparation}
\item \textbf{[C]} Both axes, case A, $b_p=5/6$, to $0.9$ of the final $P_{00}$: 32\,ms from $\bar n_0=30$.
\item \textbf{[N]} Saturation ramp, case A, $b_p=1/2$, from $\bar n_0=30$ to $P_0=0.9$ (master equation): best constant $s=2$: 15.6\,ms; two stages $s_{\max}\to1$: 5.8, 7.6, 10.2\,ms for $s_{\max}=10$, 5, 3, i.e.\ \textbf{2.7, 2.1, 1.5$\times$ faster}; the optimal control solution (chain) is within 2\% of two stages. Constant-$s$ capture scales as $s^{-1/2}$, not $1/s$. \more{d:prep}{ramp}
\end{itemize}

\section*{Mode purity}
\textbf{[N]} At fixed $s$, $P_0$ is unchanged up to an LP$_{01}$ amplitude $\epsilon\approx3\,\etaeff$; $W_0$ halves at $\epsilon\approx3.6\,\etaeff$ ($s=1$) and $2.2\,\etaeff$ ($s=3$). The $\langle n\rangle$-based $W$ looks more fragile ($1.4\,\etaeff$, case A, $s=3$) because of the tail. LP$_{01}$ power $<10^{-3}$ keeps $\ge90\%$ of $W_0$ at $s\le1$. \more{d:purity}{purity}

\section*{Rb lines}
\textbf{[N]} Case A, all excited levels, $3\Gamma$, $s=1$: D1 $F'=2$: $W_0=2.2\times10^3\,$s$^{-1}$, $P_0\approx0.954$. D1 $F'=1$, probe on $F=2$: 3.2--$3.6\times10^3\,$s$^{-1}$, $P_0\approx0.976$. D2 $F'=2$ is similar to D1 $F'=2$ at $3\Gamma$; D2 $F'=1$ is unreliable. To $P_0=0.99$: D1 $F'=1$ is $\approx2\times$ faster. \textbf{[N, one axis] Zeeman:} a fixed polarisation leaves dark states at every $n$ (2 of 5 and 1 of 3 for D1 $F'=1$), so remixing is required. Best $W_0/\nu_r$ at $P_0\ge0.97$: closed $\Lambda$ 0.077; lin$\perp$lin with repumper 0.0225; Zeeman $\Lambda$ inside $F=2$ 0.022: about $3\times$ slower. \more{d:lines}{lines}

\section*{Memory}
\textbf{[M]} $|\phi(T)|^2\to P_0^2$ at long $T$, so efficiency is capped near $P_0^4$. Limits: lattice scattering and \textbf{axial shift: 1\% of lattice intensity gives $\approx\SI{8}{Hz}$ at an antinode.} Cool deep, store shallow. LP$_{0m}$ signal: OD and bandwidth $\times g=2.33$--6.32; efficiency only at low OD; LP$_{11}$ signal on the node: useless. \more{d:memory}{memory}
\textbf{[M]} ${\rm OD}=\sigma N\kappa/A_{\rm eff}$: set by the atom number $N$ and the column--mode overlap $\kappa$, not by the column length. Efficiency 0.9 needs $N\approx1.3\times10^4$ at $a=\SI{3.5}{\upmu m}$, $2.4\times10^5$ at \SI{15}{\upmu m} (clock $\Lambda$). \more{d:memory}{memory}

\section*{Not yet done}
Zeeman model with both axes; real fibre modes (TE/TM/HE, polarisation, beating); why the 2D rate is $0.2\,W_0$; a derivation of the dip at $n\approx0.5\,n_c$; excited-state light shift at large $\nu_z$.
\end{document}
